\section{Задание}

В рамках выполнения лабораторной работы \textit{по варианту №590834} необходимо написать сценарий для утилиты \href{https://ant.apache.org/}{Apache Ant}, который реализует компиляцию, тестирование и упаковку в jar-архив кода проекта из \href{https://github.com/aaalvaaas/web3}{лабораторной работы №3} по дисциплине «Веб-программирование».

Требования к выполнению:
\begin{itemize}
  \item Каждый этап должен быть выделен в отдельный блок сценария.
  \item Все переменные и константы, используемые в сценарии, должны быть вынесены в отдельный файл параметров.
  \item \texttt{MANIFEST.MF} должен содержать информацию о версии и о запускаемом классе.
\end{itemize}

Cценарий должен реализовывать следующие цели:
\begin{enumerate}
  \item \texttt{compile} --- компиляция исходных кодов проекта.
  \item \texttt{build} --- компиляция исходных кодов проекта и их упаковка в исполняемый jar-архив. Компиляцию исходных кодов реализовать посредством вызова цели \texttt{compile}.
  \item \texttt{clean} --- удаление скомпилированных классов проекта и всех временных файлов \textit{(если они есть)}.
  \item \texttt{test} --- запуск junit-тестов проекта. Перед запуском тестов необходимо осуществить сборку проекта \textit{(цель \texttt{build})}.
  \item \texttt{native2ascii} --- преобразование \href{https://docs.oracle.com/javase/8/docs/technotes/tools/windows/native2ascii.html}{native2ascii} для копий файлов локализации \textit{(для тестирования сценария все строковые параметры необходимо вынести из классов в файлы локализации)}.
  \item \texttt{doc} --- добавление в \texttt{MANIFEST.MF} MD5 и SHA-1 файлов проекта, а также генерация и добавление в архив javadoc по всем классам проекта.
  \item \texttt{xml} --- валидация всех xml-файлов в проекте.
  \item \texttt{scp} --- перемещение собранного проекта по scp на выбранный сервер по завершению сборки. Предварительно необходимо выполнить сборку проекта \textit{(цель \texttt{build})}.
  \item \texttt{music} --- воспроизведение музыки по завершению сборки \textit{(цель \texttt{build})}.
  \item \texttt{diff} --- осуществляет проверку состояния рабочей копии, и, если изменения касаются классов, указанных в файле параметров, выполняет commit в репозиторий svn.
  \item \texttt{history} --- если проект не удаётся скомпилировать \textit{(цель \texttt{compile})}, загружается предыдущая версия из репозитория svn. Операция повторяется до тех пор, пока проект не удастся собрать, либо не будет получена самая первая ревизия. Если работающая ревизия найдена, формируется файл с результатом \textit{diff} для изменений, сломавших сборку.
  \item \texttt{alt} --- создаёт альтернативную версию программы с измененными именами переменных и классов \textit{(используя задание \texttt{replace/replaceregexp})} и упаковывает её в jar-архив через цель \texttt{build}.
  \item \texttt{report} --- в случае успешного прохождения тестов сохраняет отчет junit в формате xml, добавляет его в репозиторий svn и выполняет commit.
  \item \texttt{env} --- осуществляет сборку и запуск программы в альтернативных окружениях; окружение задается версией java и набором аргументов виртуальной машины в файле параметров.
\end{enumerate}